\homework{12}{Wed 5 Apr 2023 16:30}{}

\begin{itemize}
	\item If $h \in L^1(\Omega, \mathcal{A}, \mu)$, show that there is a sequence $g_n : \Omega \to \mathbb{R}$ of simple functions such that
		\[
			\int _\Omega (g_n - h) d \mu \to  0
		.\] 
\begin{proof}
	WLOG we may assume that $h \ge 0$, and result for more general $h$ would follow from splitting into positive and negative parts.

	Since $h$ is unsigned integrable, there exists a sequence of simple functions $g_n$ such that $0\le g_n \le h$, $g_n$ monotone increasing, and $g_n \to h$ a.e. Then 
	\[
		\int _\Omega (h - g_n) d \mu \le \int _\Omega h d \mu < \infty
	.\] 
	We know that $(h - g_n)$ is monotone decreasing. Hence by downward monotone convergence theorem,
	\[
		\int _\Omega(h - g_n) d \mu \to 0
	.\] 
	Hence
	\[
		\int _\Omega(g_n - h) d \mu \to  0
	.\] 
\end{proof}

	\item A function $g$ on an interval $[a,b]$ is piecewise constant if there are $a = t_0 < t_1 < \ldots < t_m = b$ such that the restrictions $g|_{(t_{j-1}, t_j)}$ are constant for all $j$. If $f \in L^1[a,b]$ and $\varepsilon>0$, show that there is a piecewise constant $g: [a,b] \to  \mathbb{R}$ such that 
		\[
			\int _a^b |f-g| < \varepsilon
		.\] 
\begin{proof}
	We can have $L^1$ approximation of $f$ by simple functions to arbitrary precision, so it suffices to prove the claim assuming  $f$ is a simple function. Furthermore, since $f$ is finite linear combination of indicator functions of measurable sets, and that finite linear combinations of piecewise constant functions are still piecewise constant, we may assume $f$ is an indicator function of some Lebesgue measurable $E \subset [a,b]$.

	But for any Lebesgue measurable $E \subset [a,b]$ , we can approximate $E$ by an open set $A$ such that $m(E \Delta A)$ can be arbitrarily small, where $\Delta$ is the symmetric difference. Hence we may assume $E$ is open. But then $E$ is a countable union of disjoint open cells $I_n$, and
	 \[
		m(E) = \sum_{n=1}^{\infty} |I_n| < b-a
	.\] 
	By truncating the sum at some $N$, we can have $E_N = \bigcup_{n \le N} I_n$ such that
	\[
		\int_{[a,b]} |\chi_{E} - \chi_{E_N}| = \sum_{n = N+1}^{\infty }|I_n| < \delta
	.\] 
	where $\delta$ can be arbitrarily small. Note that $\chi_{E_N}$ is piecewise constant, so we are done.
\end{proof}

	\item Suppose $\varphi_j \in L(\Omega, \mathcal{A}, \mu)$ are uniformly bounded, $j \in \mathbb{N}$, and $\sup  \left\{ \left| \varphi_j(x) \right| : j \in \mathbb{N} , x \in \Omega\right\}  <\infty $. Also assume $\mu(\Omega) < \infty $. If $\int _\Omega \varphi_i \varphi_j = 0$ whenever $i \neq j$, prove that the average 
		\[
			\psi_n = \frac{\varphi_1 + \varphi_2 + \ldots + \varphi_{n^2}}{n^2} \to 0 \qquad \text{a.e.}
		.\] 
\begin{proof}
	Note that
	\[
		\int_\Omega \psi_n^2 = \frac{1}{n^4}\int _\Omega( \varphi_1 ^2 + \ldots + \varphi_{n^2}^2) \le \frac{1}{n^2} M^2 \mu(\Omega)
	.\] 
	Thus following series converges:
	\[
		\sum_{n = 0}^\infty  \int_\Omega \psi_n^2
	.\] 
	By monotone convergence theorem, we may exchange sum and integral:
\[
		\int_\Omega \sum_{n = 0}^\infty  \psi_n^2< \infty 
	.\]
	Hence $\sum_{n = 0}^\infty  \psi_n^2$ is finite a.e., Hence  $\psi_n \to 0$ a.e.
\end{proof}

	\item If $f, g: [a,b] \to \mathbb{R}$ are absolutely continuous, prove that so is $fg$.
\begin{proof}
	Let $\delta>0$, and choose $\delta>0$ be as in the definition of absolutely continuity of $f$. Assume that the same $\delta$ also satisfies absolute continuity of $g$.

	$f,g$ absolutely continuous implies $f, g$ are of bounded variation, implies $f, g$ bounded on $[a,b]$. Let $M$ be the bound.

	Let  $I_j$ be disjoint intervals in  $[a,b]$ with total length less than $\delta$. We want to bound the quantity
	\[
		\sum_{j = 1}^n [fg]_{I_j}
	.\] 
	Note that if $I = [\alpha, \beta]$,
	\[
		[fg]_{I} = |f(\beta) [g]_{I} + g(\alpha) [f]_{I}| \le M \left( [f]_I + [g]_I \right) 
	.\] 
	Hence
	\[
		\sum_j [fg]_{I_j} \le  M \left( \sum_j [f]_{I_j} + \sum_j [g]_{I_j} \right) 
	.\] 
	By using absolute continuity of $f$ and $g$, we can get
	\[
		\sum_j [fg]_{I_j} \le  2 M \varepsilon
	.\]
	Hence, by choosing small enough $\delta$, we can make sum $\sum_j [fg]_{I_j}$ arbitrarily small, where $I_j$ is any finite collection of intervals whose total length is smaller than $\delta$. Hence $fg$ is absolutely continuous.
\end{proof}

	\item  Prove that any absolutely continuous function $\varphi:[a,b] \to \mathbb{R}$ is of bounded variation.
\begin{proof}
	Fix $\varepsilon>0$ and  choose $\delta>0$ as in the definition of absolutely continuity of $\varphi$. Let $P$ be a partition of $[a,b]$ where each cell has length at most $\delta$. We focus on one of the cells, call it $A = [\alpha, \beta]$. 

	For any partition of $A$, say $Q: x_0 < \ldots < x_n$, the partition can be viewed as finite union of disjoint intervals with total length less than $\delta$. Thus
	\[
		\sum_{i = 1}^n |F(x_i) - F(x_{i-1})| < \varepsilon
	.\] 
	If we take supremum over the partitions, we obtain
	\[
		\|\varphi\|_{TV(A)} \le \varepsilon
	.\] 
	By property of $\varphi$,
	\[
		\|\varphi\|_{TV([a,b])} = \sum_{A \in P}\|\varphi\|_{TV(A)} \le N \varepsilon
	.\] 
	where $N$ is the number of intervals in the partition $P$. We see that $N \varepsilon$ is finite, so $\varphi$ has bounded variation.
\end{proof}

	\item If $g:[a,b] \to \mathbb{R}$ is increasing and absolutely continuous, prove that the associated Lebesgue-Stieltjes measure $m_g$ is absolutely continuous with respect to Lebesgue measure $m$, i.e.  $m_g \ll m$.
\begin{proof}
	Take $E \subset [a,b]$ such that $m(E) = 0$. We want to prove $m_g(E) = 0$. Fix $\varepsilon>0$, we want to cover $E$ with disjoint open intervals $I_n$ where $\sum_{n = 1}^\infty  m_g(I_n) < \varepsilon$. 

	Take $\delta$ as in the definition of absolute continuity of $g$.
	Since $m(E) = 0$, we can pick disjoint open intervals $I_n$ where $\sum_{n = 1}^\infty  m(I_n) < \delta$. Then by absolute continuity,
	\[
		\sum_{n = 1}^{\infty } [g]_{I_n} < \varepsilon
	.\] 
	By definition of $m_g$, for $I_n = (a, b)$,
	\[
		m_g (I_n) = \lim_{b^-} g - \lim_{a^+} g = g(b) - g(a) = [g]_{I_n}
	.\] 
	Hence $\sum_{n = 1}^\infty  m_g(I_n) \le \sum_{n = 1}^\infty  [g]_{I_n} < \varepsilon$.
\end{proof}

	\item If $f, g$ are absolutely continuous on $[a,b]$, prove $\int _a^b f' g = f g|_a^b - \int _a^b fg'$.
\begin{proof}
	Since $f, g$ are absolutely continuous, so if $fg$. Hence $fg$ is differentiable a.e. and
	\[
		\int _a^b (fg)' = fg|_a^b
	.\] 
	By product rule of differentiation, $(fg)' = f'g + fg'$, apply linearity of integral and we are done.
\end{proof}

	\item If $g:[a,b] \to \mathbb{R}$ is increasing and absolutely continuous, and $m_g$ is the associated Lebesgue-Stieltjes measure, prove that for any Borel set $E \subset [a,b]$, we have
		\[
			\int _E g' = m_g (E)
		.\] 
\begin{proof}
	When $E$ is an open interval $(a,b)$, using definitions and continuity of $g$,
	\[
		\int _E g' = g|_a^b = g(b) - g(a) = m_g(E)
	.\] 
	Every open set is a countable union of disjoint open intervals. We can use upward MCT on both sides of the equation, to extend the result to open sets.

	Each closed set is countable intersection of open sets. By applying intersections one at a time, we arrive at a decreasing sequence of open sets that converge to that closed set. Apply downward MCT to both sides of the equation, we can extend the result to closed sets.

	Every Borel set can be approximated from inside by a compact set $B$ and from above by an open set $A$. Pick those two sets close enough to each other such that $m_g(A \setminus B) < \varepsilon$. Then we have both
	\[
		\int _B g' \le \int _E g' \le \int _A g'
	.\] 
	since $g' \ge 0$, and
	\[
		m_g(B) \le m_g(E) \le m_g(A)
	.\] 
	by monotonicity of measure.

	By what we have proved, $\int _B g' = m_g(B)$ and $\int _A g' = m_g (A)$. Also $m_g(A) - m_g(B) < \varepsilon$, so
	\[
		\left| \int _E g' - m_g(E) \right| <  \varepsilon
	.\] 
	Since $\varepsilon$ is arbitrary, 
	\[
		\int _E g'  = m_g (E)
	.\] 
	for arbitrary Borel set $E$.
	

\end{proof}

\end{itemize}

